\documentclass[11pt,a4paper]{article}
\usepackage[T1]{fontenc}
\usepackage[utf8]{inputenc}
\usepackage[margin=1in]{geometry}
\usepackage{amsmath,amssymb,amsthm}
\usepackage{stmaryrd}
\usepackage{booktabs}
\usepackage{tabularx}
\usepackage{array}
\usepackage{enumitem}
\usepackage{listings}
\usepackage{xcolor}
\usepackage[expansion=false]{microtype}
\usepackage{tikz}
\usetikzlibrary{arrows.meta,positioning,fit,calc}
\usepackage[htt]{hyphenat}
\usepackage[hidelinks,pdftitle={F1R3Comb-Mat v0.4}]{hyperref}
\setlength{\emergencystretch}{3em}

\newcommand{\code}[1]{\texttt{#1}}
\newcommand{\Comb}{\textsc{F1R3Comb}}
\newcommand{\Mat}{\textsc{F1R3Comb-Mat}}
\newcommand{\Kind}{\textsc{K1ndl1ng}}
\newcommand{\Camp}{\textsc{CampF1R3}}
\newcommand{\Bell}{\textsc{B3ll0ws}}
\newcommand{\Fx}{\textsc{F1R3x}}
\newcommand{\Rhobots}{\textsc{Rhobots}}
\newcommand{\MUST}{\textsc{must}}
\newcommand{\MUSTNOT}{\textsc{must not}}
\newcommand{\SHOULD}{\textsc{should}}
\newcommand{\SHOULDNOT}{\textsc{should not}}
\newcommand{\MAY}{\textsc{may}}

\newcommand{\quo}[1]{@#1}
\newcommand{\drp}[1]{{*}#1}
\newcommand{\nil}{\mathbf{0}}
\newcommand{\nequiv}{\equiv_{\mathcal{N}}}
\newcommand{\nms}[1]{\mathcal{N}^{*}(#1)}
\newcommand{\Nsort}{\mathsf{N}}
\newcommand{\Psort}{\mathsf{P}}
\newcommand{\red}{\longrightarrow}
\newcommand{\wred}{\longrightarrow^{*}}
\newcommand{\scong}{\equiv}
\newcommand{\comps}[1]{\lvert #1 \rvert}

\newcommand{\mm}{\mathsf{m}}
\newcommand{\dd}{\mathsf{d}}
\newcommand{\kk}{\mathsf{k}}
\newcommand{\fw}{\mathsf{fw}}
\newcommand{\bl}{\mathsf{b}_{\mathsf{l}}}
\newcommand{\br}{\mathsf{b}_{\mathsf{r}}}
\newcommand{\sy}{\mathsf{s}}
\newcommand{\ev}{\mathsf{e}}
\newcommand{\qq}{\mathsf{q}}
\newcommand{\cons}{\mathsf{cons}}
\newcommand{\consm}{\mathsf{cons}_{\mm}}
\newcommand{\consd}{\mathsf{cons}_{\dd}}
\newcommand{\conss}{\mathsf{cons}_{\sy}}
\newcommand{\conse}{\mathsf{cons}_{\ev}}
\newcommand{\consf}{\mathsf{cons}_{\fw}}

\newcommand{\Spec}{\mathcal{S}}
\newcommand{\Nm}{\mathcal{N}}
\newcommand{\cre}[1]{a^{\dagger}_{#1}}
\newcommand{\ann}[1]{a_{#1}}
\newcommand{\Wt}{\mathsf{W}}

\theoremstyle{definition}
\newtheorem{definition}{Definition}[section]
\newtheorem{requirement}[definition]{Requirement}
\newtheorem{decision}[definition]{Decision}
\newtheorem{obligation}[definition]{Obligation}
\newtheorem{hazard}[definition]{Hazard}
\newtheorem{erratum}[definition]{Erratum}
\newtheorem{remark}[definition]{Remark}
\newtheorem{example}[definition]{Example}
\theoremstyle{plain}
\newtheorem{proposition}[definition]{Proposition}
\newtheorem{lemma}[definition]{Lemma}

\newcommand{\Lft}{\mathsf{L}}
\newcommand{\Rgt}{\mathsf{R}}
\newcommand{\inst}{\mathsf{inst}}
\newcommand{\ERECT}{\mathrm{ERECT}}
\newcommand{\ALLOC}{\mathrm{ALLOC}}
\newcommand{\enc}[1]{\llbracket #1 \rrbracket}
\newcommand{\inp}[3]{\mathsf{for}(\,#1 \leftarrow #2\,)#3}
\newcommand{\out}[2]{#1!(#2)}
\newcommand{\Fam}{\mathcal{F}}

\lstset{
  basicstyle=\ttfamily\footnotesize,
  columns=fullflexible,
  keepspaces=true,
  breaklines=true,
  frame=leftline,
  framesep=6pt,
  xleftmargin=8pt,
  rulecolor=\color{black!30},
  showstringspaces=false,
}

\newcommand{\cstar}{\mathsf{cons}^{*}}

\title{\textbf{\Mat}\\[1mm]
\large The rho combinators as sparse linear algebra:\\
a Rust implementation of the matrix representation,\\
the bulk-synchronous machine, and the transfer operator\\[2mm]
\normalsize Specification v0.4 --- companion to \Comb{} v0.5}
\author{F1R3FLY.io}
\date{23 September 2026}

\begin{document}
\maketitle

\begin{abstract}
\noindent
\Mat{} is the third stage of the route from the rho calculus to a
data-parallel device: it takes a \Comb{} term to the matrix representation of
\cite{matnote} and runs it on a bulk-synchronous machine whose step is a
maximal independent set of redexes. Phase two of the translation erects each
run-time instance at an address, and draft~3 shows that doing so within the
discipline its correctness proof needs requires a single-premise rule that
places the address more than once \cite[Prop.~6.17, Lem.~6.19]{draft3}. It
offers two such rules \cite[\S6.7]{draft3}: curried path-indexed
constructors, and the context-labelled combinator
$\inst(a,f,C) \mid \mm(a,v) \red \mm(f,\quo{C[v]})$.

The implementation path is built on $\inst$, and this version makes that
explicit. $\inst$ is the machine's primary erection: the default build, the
only scheme the device contract targets, and the one against which width,
cost and conformance are measured. The evidence of \code{phase2.py} supports
the choice on every axis the machine cares about. Both schemes erect
top-level atoms without mixing (E5). On the recursion combinator $\inst$ takes
7 parallel steps and builds 16 names per unfolding against the curried
scheme's 10 and 22 (E6). And $\inst$ reaches every gate store, while the
curried scheme can assemble a store whose body mentions scoped names in two
atoms only by a join, which mixes 44\% of stores at two instances and 87\% at
eight (E7). $\inst$ also keeps every row four names wide, has the cheapest
join available, and erects a unit in two steps regardless of its size, so
erection adds width rather than depth.

\Comb{} v0.5 currently defaults to the curried scheme
\cite[Dec.~5.19]{combspec}. \Mat{} runs curried artefacts on the host so that
it can execute what \Comb{} emits today, but the device path does not carry
them, and the first request this version puts to \Comb{} is to make $\inst$
its default, with the reasons v0.5 gave for the opposite choice answered
from the machine's side. The rest of v0.3 stands: one instantiation kernel
with a validity class for curried templates, the literal erection as a
negative control, the marking read from the header, provenance, the marking
oracle.
Key words \MUST, \MUSTNOT, \SHOULD, \MAY{} are used in the sense of RFC~2119.
\end{abstract}

\tableofcontents

%======================================================================
\section{Purpose and scope}\label{sec:purpose}

\subsection{The route, and where this document sits on it}

\begin{center}
\begin{tabularx}{\linewidth}{@{}l>{\raggedright\arraybackslash}Xl@{}}
\toprule
Stage & What it does & Specified in\\
\midrule
1 & K0 text to a canonical normal form & \cite{k1ndl1ng}\\
2 & Normal form to combinator term, in two phases & \Comb{} v0.5 \cite{combspec}\\
3 & Combinator term to matrix representation; the bulk-synchronous machine;
    the transfer operator & this document\\
4 & The machine's kernels on a device & Section~\ref{sec:device}, as a
    contract\\
\bottomrule
\end{tabularx}
\end{center}

\noindent
This version fixes $\inst$ as the machine's primary erection. Otherwise it
leaves v0.3's changes to phase two, and the rest of v0.2 \cite{matv02}, as
they were. Where a requirement below says ``as
v0.2'', v0.2's text is normative; where the two differ, this document
governs. Two properties govern everything below. The machine computes exactly one step
per state, seed and configuration, so a host and a device can be compared
trace for trace (\cite[P4]{matnote}); and it runs what stage two emits,
including every instance of every unit, at the width the program has
(\cite[Princ.~10.4]{matnote}).

\subsection{What exists today}

\begin{center}
\begin{tabularx}{\linewidth}{@{}l>{\raggedright\arraybackslash}X@{}}
\toprule
Artefact & What it gives this project\\
\midrule
\code{rho-combinators-draft3} & The calculus and its presentations
(\S3.2--3.3); the constructor family (Def.~3.8, Rem.~3.9); erection at a
computed name (Rem.~3.11); phase one (\S6.2) and phase two (\S6.3, Def.~6.8,
Ex.~6.9); the discipline and the permutation lemma (Def.~6.10, Lem.~6.11);
freshness by size (Lem.~6.14); \emph{new:} the literal erection's failure
(Ex.~6.16), multiplicity (Prop.~6.17, Def.~6.18, Lem.~6.19), the marking
(Rem.~6.20), the deciding term (Ex.~6.21), the two schemes and the
recommendation (\S6.7, Rem.~6.22); the weighted count (Lem.~9.2)\\
\code{f1r3comb-spec-v05} & The term type and shape table; tags; phase two as
a compiler; \emph{new:} the literal erection forbidden (Req.~5.15), v0.4's
two-stage erection withdrawn (Rem.~5.16), the marking analysis (Req.~5.17),
the deciding term as a test (Req.~5.18), curried by default
(Dec.~5.19--Req.~5.23), the mixing and marking obligations (Obl.~11.3,
11.4)\\
\code{rhocomb-matrix} v1.2 & The state space, incidence, step and operator;
\S14, phase two seen from the machine\\
\code{rhocombmat.py} v1.3, \code{phase2.py}, \code{phase2-results.txt} & The
reference machine with the generated table, $\inst$, and the curried
constructor $\cstar$; the three-arm experiments E5--E7 of
Section~\ref{sec:evidence}\\
\code{rhocomb-logic} rev.~2 & Components (Def.~2.9); drop elimination
(Lem.~2.8); injective quotation (Prop.~2.10); the weighted count as a formula
(Def.~4.2, Prop.~6.17)\\
\bottomrule
\end{tabularx}
\end{center}

\subsection{Nomenclature}\label{sec:nomen}

\begin{center}
\begin{tabularx}{\linewidth}{@{}l>{\raggedright\arraybackslash}X@{}}
\toprule
Term & Meaning in this document\\
\midrule
\emph{term}, \emph{shape table}, \emph{family} & As in v0.2: a closed
drop-free \code{Term}; \Comb's table of shapes; the constructors a program
uses, $\Fam(P)$\\
\emph{row width}, $w$ & The largest arity in the program's signature\\
\emph{name id}, \emph{row}, \emph{row ref} & As in v0.2\\
\emph{address}, \emph{leaf}, \emph{static channel}, \emph{unit},
\emph{instance} & As in v0.2 and \cite[\S1.3]{combspec}\\
\emph{scoped}, \emph{marked} & A bound name allocated beneath the address; a
name carrying an instance-specific value, in the sense of
\cite[Rem.~6.20]{draft3}\\
\emph{template} & An interned name over the reserved \code{HOLE}: the thing
an erecting rule instantiates at an address\\
\emph{scheme} & How a unit is erected: \emph{curried}
\cite[\S6.7]{draft3}, \emph{$\inst$} \cite[\S6.7]{draft3}, or \emph{literal}
\cite[Def.~6.8]{draft3}\\
\emph{provenance} & The instance a row belongs to, as far as the machine can
tell; metadata, never semantics\\
\emph{mixed} & An atom mentioning scoped names of two addresses\\
\bottomrule
\end{tabularx}
\end{center}

\subsection{Goals}

\begin{enumerate}[leftmargin=1.6em]
\item A total, lossless, round-trippable map from \Comb{} terms to the matrix
representation, for every shape \Comb's shape table generates and every
scheme \Comb{} emits.
\item Redex finding as joins derived from the generated rule table, with three
finders that agree.
\item A bulk-synchronous step defined bit-exactly.
\item Phase-two images erected by the context-labelled combinator $\inst$
run correctly, with \Comb's allocator, freshness checks and marking; curried
artefacts run on the host for compatibility with \Comb{} v0.5; the literal
scheme is available as a negative control.
\item The instruments that bear on \Comb's open choices: width with erection
separated, per-scheme cost, cross-instance mixing, and a brute-force oracle
for the marking analysis.
\item On the destructor-free sublattice, the transfer operator, scoped
honestly.
\item \Comb's portability envelope.
\end{enumerate}

\noindent
Non-goals are those of v0.2: device code (a contract only), presentation~B,
weights, the logic, the history unfolding, intern-table collection, and phase
one and two themselves. \Mat{} takes a \code{Term} and \MUSTNOT{} depend on
\code{f1r3comb-lower} or \code{f1r3comb-ir}.

%======================================================================
\section{Position relative to \Comb}\label{sec:position}

\begin{center}
\begin{tabularx}{\linewidth}{@{}l>{\raggedright\arraybackslash}X>{\raggedright\arraybackslash}X@{}}
\toprule
& \Comb{} machine (\code{f1r3comb-run}) & \Mat{} (\code{f1r3comb-par})\\
\midrule
State & \Camp{} store keyed by \code{Hash32} & shape tables of \code{NameId}
rows, width $w$\\
Name identity & content hash of the encoding & structural intern in-run,
content hash at the boundary\\
Event & one comm & one step, a maximal independent set\\
Resolver & store order key; adversarial profile \cite[Obl.~11.2]{combspec}
& priority function; cross-instance profile\\
Rules & generated from the shape table & generated from the same table\\
Erection & curried by default; $\inst$ under \code{inst-erection} &
$\inst$ primary; curried on the host for compatibility; literal as a negative
control\\
Addresses & allocator \cite[Req.~9.5]{combspec} & the same allocator\\
\bottomrule
\end{tabularx}
\end{center}

\subsection{Decisions}

\begin{decision}[Presentation A only]\label{dec:aonly}
Unchanged from v0.2. Under B rows are ragged and (G1) of
\cite[\S2.1]{matnote} fails; the A/B comparison \cite[Obl.~11.10]{combspec}
is run on \code{f1r3comb-run}.
\end{decision}

\begin{decision}[A store row holds the quotation of its payload]
\label{dec:qrow}
Unchanged from v0.2: the row for $\qq(a,p)$ is
$(\qq,\mathrm{id}(a),\mathrm{id}(\quo{p}))$, lossless under A, with the shape
tag carrying the $\mm$/$\qq$ distinction, and $\cons_{\qq}$ uniform with every
other member.
\end{decision}

\begin{decision}[Decode by arena read]\label{dec:decode}
Unchanged from v0.2, and still at odds with \cite[Req.~9.3]{combspec}
(Section~\ref{sec:paper}, item~9).
\end{decision}

\begin{decision}[The rule table is generated from \Comb's shape table]
\label{dec:ruletable}
Unchanged from v0.2. The generator lives in \code{f1r3comb-term}; \Mat{} lists
no rules of its own.
\end{decision}

\begin{decision}[$\inst$ is the primary erection]\label{dec:schemes}
The implementation path erects by the context-labelled combinator $\inst$.
\Mat{}'s default feature set \MUST{} enable \code{inst-erection}, the feature
name \Comb{} uses, so that one build configuration means the same thing on
both machines. Every measurement that bears on the path --- width, cost per
instance, conformance, the device conformance corpus --- \MUST{} be made on
$\inst$ artefacts first, and the device contract targets $\inst$ alone
(Requirement~\ref{req:device}). \Mat{} runs whatever scheme the artefact
header names, and \MUST{} also run curried artefacts on the host, because
curried is \Comb's default today \cite[Dec.~5.19]{combspec}; that support is
for compatibility and differential testing, not a second supported path, and
is asked to be retired by Section~\ref{sec:paper}, item~1. It \MUST{} support the
literal scheme under a test-only feature \code{literal-erection}, which no
product crate enables, because \cite[Obl.~11.3]{combspec} requires the broken
construction as a negative control and \Comb{} is forbidden to emit it
\cite[Req.~5.15]{combspec}. It \MUST{} accept units that \Comb's marking
analysis has compiled with the fixed-arity family \cite[Req.~5.17]{combspec},
reading the marking from the header (Section~\ref{sec:marking}).
\end{decision}

\begin{decision}[One kernel, whose primary client is $\inst$]
\label{dec:onekernel}
A curried constructor $\cstar_{A}[\vec w](t,f)$ and an $\inst(a,f,C)$ both
consume one message and deliver the quotation of a template instantiated at
its payload. \Mat{} \MUST{} implement both by one rule shape with one encode,
\code{Instantiate\{template, with\}}, and distinguish them by a validity class
on the template (Requirement~\ref{req:template}). \cite[Rem.~5.22]{combspec}'s
second reason for the curried default --- that $\inst$ forfeits the property
that makes \Comb's runtime a store with a three-line matcher --- is a property
of \code{f1r3comb-run}; on this machine both schemes are the same kernel, and
the curried scheme's leaves are built by the same substitution
(Section~\ref{sec:paper}, item~4).
\end{decision}

%======================================================================
\section{Phase two, with the evidence}\label{sec:evidence}

\subsection{The obstruction, now in the paper}

\cite[\S6.6]{draft3} shows the literal erection fails the discipline on a
single instance (Ex.~6.16), that the obvious repair regresses, and that the
obstruction is structural: under the discipline extended to all rules every
reachable atom holds at most one occurrence of an address
(Prop.~6.17), so erecting any atom with two scoped names needs a rule of
multiplicity at least two (Lem.~6.19). v0.2's Section~3.2 is discharged by
this and \Comb{} v0.5 withdraws its two-stage construction
\cite[Rem.~5.16]{combspec}. \cite[\S6.7]{draft3} supplies two single-premise
rules of the required multiplicity:
\[
  \cstar_{A}[\vec w](t,f) \mid \mm(t,\sigma) \red
  \mm(f,\quo{A(\sigma\!\cdot\!w_{1},\dots,\sigma\!\cdot\!w_{n})}),
  \qquad
  \inst(a,f,C) \mid \mm(a,v) \red \mm(f,\quo{C[v]}).
\]

\subsection{Three arms}

\cite[Rem.~5.20]{combspec} observed that v0.2's measurements compared the
literal scheme with $\inst$ and had no curried arm. \code{phase2.py} now runs
all three, with the curried scheme implemented as \cite[Rem.~5.21]{combspec}
describes it: no allocation gadget, one $\dd$ per erected atom fanning the
address out on static channels, one $\cstar$ and one $\ev$ per atom. Its
templates allow a static name at an argument position, as
\cite[Ex.~6.16]{draft3}'s $\dd(a,p_{0},q_{1})$ requires.

\begin{center}\small
\begin{tabular}{@{}lrrrr@{}}
\toprule
\multicolumn{5}{@{}l}{\textbf{E5.} Unit $\{\fw(z_{1},z_{2}),\ \kk(z_{1})\}$;
runs with a mixed atom, of 50}\\
\midrule
instances & 2 & 4 & 8 & 16\\
\midrule
literal, maximal progress / sequential & 29 / 16 & 48 / 38 & 50 / 49 & 50 / 50\\
curried, either schedule & 0 & 0 & 0 & 0\\
$\inst$, either schedule & 0 & 0 & 0 & 0\\
\midrule
parallel steps to erect: literal 5, curried 3, $\inst$ 2\\
\bottomrule
\end{tabular}
\end{center}

\begin{center}\small
\begin{tabular}{@{}lrrr@{}}
\toprule
\multicolumn{4}{@{}l}{\textbf{E6.} $D_{x}$ of \cite[Ex.~6.9]{draft3}, 200
unfoldings, per unfolding}\\
\midrule
scheme & parallel steps & sequential steps & new names\\
\midrule
curried & 10 & 32.0 (29--35) & 22\\
$\inst$ & 7 & 10.0 (9--11) & 16\\
\bottomrule
\end{tabular}
\end{center}

\begin{center}\small
\begin{tabular}{@{}lrr@{}}
\toprule
\multicolumn{3}{@{}l}{\textbf{E7.} Store $\qq(b,\quo{(\fw(c_{1},c_{3}) \mid
\kk(c_{2}))})$, $b$ static, $c_{i}$ scoped; mixed stores}\\
\midrule
scheme & 2 instances & 8 instances\\
\midrule
curried, payload joined by $\cons_{\mid}$ & 44 / 100 & 347 / 400\\
$\inst$ & 0 / 100 & 0 / 400\\
\bottomrule
\end{tabular}
\end{center}

Three conclusions, each a requirement below.

\emph{Both conforming schemes are correct for top-level atoms.} E5 settles
\cite[Rem.~5.20]{combspec}'s question on conformance: no mixing under either,
at any count, under either schedule. Conformance does not separate them;
cost and reach do.

The three together are why the implementation path uses $\inst$.

\emph{On this machine the curried scheme does more work, not less.}
\cite[Rem.~5.22]{combspec}'s first reason is that $\inst$'s two steps hide
$O(|\mathrm{unit}|)$ work in one rule, recomputing encodings and content
hashes up the rebuilt term. Structural interning computes no hashes
(Requirement~\ref{req:hashfree}), so the work on this machine is the names
interned: 16 per unfolding under $\inst$ against 22 under curried, where each
$\cstar$ builds its leaves by $d$ interns from the address and the duplicator
tree adds depth. Both schemes should be charged by names built, not by steps,
and on that measure $\inst$ is cheaper here.

\emph{The curried scheme does not reach every store.} Its template places the
address at argument positions of one atom. A gate store whose body mentions
scoped names in two atoms --- which is the body of any input whose bound name
occurs in two atoms of its continuation --- needs its payload assembled from
two address-bearing pieces, and the only rule that assembles them is a join.
The store atom then mentions two marked names, which is
\cite[Rem.~6.20]{draft3}'s own criterion, so Lemma~6.19 applies to it; the
multiplicity it needs is beneath a quotation, which the curried rule does not
supply. \cite[Req.~4.9]{combspec}'s recipe (build the payload by $\cons_{\ev}$,
feed it to $\cons_{\qq}$) conforms exactly when the store's subject is static
and the payload has one scoped atom, which is the case of
\cite[Ex.~6.9]{draft3}; E6 uses it there. E7 is the case beyond it.

\begin{hazard}[Curried stores]\label{haz:curriedstores}
Under the curried scheme, a unit whose gate store's payload mentions scoped
names in more than one atom, or beneath a further quotation, cannot be erected
within the discipline, and a compiler that falls back to a join produces the
mixing E7 measures. Whether that mixing is observable is the deciding-term
question of \cite[Ex.~6.21]{draft3} applied to stores, which the marking
analysis does not yet ask (Section~\ref{sec:paper}, item~2). \Mat{} \MUST{}
detect such stores at deploy under the curried scheme and report them as
\code{comb-mat-curried-store}, a warning, naming the unit.
\end{hazard}

\subsection{Templates}

\begin{requirement}[One representation for both schemes]\label{req:template}
A template \MUST{} be an interned name over the reserved \code{NameId::HOLE}.
Both erecting rules \MUST{} be rows of arity three --- $(t,f,T)$ for $\cstar$,
$(a,f,C)$ for $\inst$ --- with the template in the third column, so the row
width is four under either scheme. A $\cstar$ template \MUST{} satisfy the
curried validity class: it is a single atom, and each argument is either a
path beneath \code{HOLE} (a $\Lft/\Rgt$ spine) or a name that does not mention
\code{HOLE}. The check runs at import and at every firing in debug builds;
a violation is \code{comb-mat-template}. An $\inst$ template is any name over
\code{HOLE}. \cite[Req.~5.23]{combspec}'s ``context node, with holes numbered
canonically'' maps onto this representation with every hole identified with
\code{HOLE}, since a single-input context has one hole variable.
\end{requirement}

\begin{requirement}[The hole]\label{req:hole}
As v0.2: \code{HOLE} has no counterpart in the term type outside a template
position, is rejected by import anywhere else (\code{comb-mat-hole}), and no
rule's conclusion may contain it. Its tag is \Comb's to allocate
(Section~\ref{sec:paper}, item~8).
\end{requirement}

\begin{requirement}[Instantiation]\label{req:subst}
The encode \MUST{} compute $T\{v/\code{HOLE}\}$ over the interned form,
hereditarily beneath quotation, with a per-firing memo keyed by name id, so
its cost is the number of distinct names in $T$. Result names are interned in
canonical order (Requirement~\ref{req:insertorder}). The report \MUST{}
attribute the names built to the scheme, which is the work figure of E6.
\end{requirement}

\begin{requirement}[The measured costs are regressions, $\inst$ first]
\label{req:dxlaw}
\code{f1r3comb-check} \MUST{} reproduce E6 as laws over the number of
unfoldings to $10^{4}$: 7 parallel steps and 16 names per unfolding under
$\inst$, with a sequential mean of 10, as the primary regression; and 10, 22
and 32 under curried as the compatibility regression. A golden
file \MUSTNOT{} be compared across schemes \cite[Dec.~5.19]{combspec}.
\end{requirement}

\subsection{Marking, and when a crossing is sanctioned}\label{sec:marking}

\cite[Req.~5.17]{combspec} compiles a unit with the fixed-arity family when no
atom of its phase-one image mentions two marked names, since crossings there
permute which instance handles which value and every answer is right. On such
a unit this machine will observe crossings --- that is what the analysis
licenses --- and must not report them as defects.

\begin{requirement}[The header's marking is read, and checked]
\label{req:marking}
\Mat{} \MUST{} read the marking from the artefact header and carry, for each
scoped name, whether it is marked. A row mentioning scoped names of two
addresses is \emph{mixed}; it is \emph{sanctioned} if at most one of those
names is marked and \emph{unsanctioned} otherwise. Unsanctioned mixed rows
\MUST{} raise a trace event at every trace level and be counted in the run
report; under a conforming scheme with a correct marking there are none, and
the test suite \MUST{} assert it under \code{CrossInstance}.
\end{requirement}

\begin{requirement}[The oracle for the marking analysis]\label{req:oracle}
\cite[Obl.~11.4]{combspec} checks the marking against exhaustive scheduling
at two and three instances. \code{f1r3comb-op}'s bounded exploration
(Section~\ref{sec:op}) \MUST{} serve as that oracle: given a unit, an instance
count and a depth bound, it explores every interleaving, collects the encoded
observations at quiescent markings, and reports whether they differ. It
\MUST{} classify $R = \inp{y}{a}{\out{y}{\drp{y}}}$ as unsafe and
$\inp{y}{a}{(\drp{y} \mid \drp{y})}$ as safe \cite[Req.~5.18]{combspec},
and \SHOULD{} also be run on E7's store, which neither \Comb{} nor the draft
has yet classified.
\end{requirement}

\subsection{Provenance}\label{sec:prov}

\begin{definition}[Provenance]\label{def:prov}
As v0.2: every row carries \code{Prov}, \code{STATIC} on import; a produced
row takes its premises' common non-static tag, \code{STATIC} if all are
static, \code{MIXED} otherwise; and a row naming a leaf takes that leaf's
root's tag from the intern table's \code{root} column. For a template
instantiation the tag is the consumed address's.
\end{definition}

\begin{requirement}[Provenance is metadata]\label{req:provmeta}
As v0.2: it \MUSTNOT{} affect redexes, selection under
\code{MaximalProgress}, or conclusions; it appears in traces as the root
address's content hash. \code{MIXED} rows are classified sanctioned or
unsanctioned by Requirement~\ref{req:marking}.
\end{requirement}

%======================================================================
\section{The representation}\label{sec:rep}

v0.2's representation is unchanged: structural interning with no hash on the
hot path; \code{size}, \code{root} and \code{depth} columns; ids never leaving
the process; the budget, including address depth; structure-of-arrays shape
tables with a provenance column; import, export and round trip; the
canonical \code{.cmat}. Three points are restated.

\begin{requirement}[Structural interning]\label{req:hashfree}
As v0.2, with the observation of E6 added: prefix sharing is by construction,
and per-unfolding name cost is constant under both schemes (16 and 22).
\end{requirement}

\begin{requirement}[Width]\label{req:soa}
$w$ is read from the header and instantiated for $w \in \{4,5,6\}$. Under the
curried scheme and under $\inst$, $w = 4$: base shapes have arity at most
three, $\cons_{\dd}$ and $\cons_{\sy}$ four, the erecting rows three. $w = 5$
arises only under \code{literal-erection} with an (o5) occurrence, which is a
negative control.
\end{requirement}

\begin{requirement}[\code{.cmat}]\label{req:cmat}
As v0.2, with the header extended by the scheme, the marking, and the
template names; \code{HOLE} relabelled to index $0$.
\end{requirement}

%======================================================================
\section{The rule table}\label{sec:rules}

\begin{requirement}[What the generator emits]\label{req:gen}
As v0.2 --- the seven routing and destructor rules, $\mathsf{quote}$,
$\mathsf{make}_{\mid}$, one rule per member of $\Fam(P)$ --- and, per the
header's scheme, the erecting rule: under curried, $\cstar$ with one loud
premise at slot 0 and conclusion
\code{Encode\{deliver: Slot(1), node: Instantiate\{template: Slot(2), with:
Payload(0)\}\}}; under $\inst$ the same with its own shape; under literal,
nothing further, the family doing the erecting. Classes, join plans, the
grading and $w$ are derived from the emitted table.
\end{requirement}

\begin{requirement}[Only a forwarder reads a store]\label{req:quietjoin}
As v0.2, a property of the generator.
\end{requirement}

\begin{requirement}[Classes and plans are checked]\label{req:plans}
As v0.2, per scheme: one class-(c) rule; one class-(b) rule per member of
$\Fam(P)$, plus $\mathsf{make}_{\mid}$ and the scheme's erecting rule; seven
class-(a) rules.
\end{requirement}

\begin{requirement}[Grading]\label{req:grading}
As v0.2. The erecting rules, like the family, have weight one.
\end{requirement}

%======================================================================
\section{Redex finding, selection, firing}\label{sec:machine}

The sparse layer, the three finders, the canonical key, the distinctness and
symmetry masks, $n$-way joins over the generated signature, the conflict
hypergraph, the \code{mix64} priority, the greedy maximal independent set and
its equality with parallel rounds, the four resolvers including
\code{CrossInstance}, race visibility, the order of effects, and debug-build
soundness checking are unchanged from v0.2.

\begin{requirement}[Order of effects]\label{req:insertorder}
As v0.2: sort the step by canonical key; compute conclusions in that order,
interning new names --- including those an instantiation creates, in a fixed
traversal of the template --- in that order; remove consumed rows by stable
compaction; append produced rows in the order computed. Parallel backends
\MUST{} assign ids as if sequentially.
\end{requirement}

\begin{hazard}[Contention at statics, per scheme]\label{haz:statics}
Under $\inst$, $n$ instances contend at one static channel $r$ per unit with a
one-premise rule: $n^{2}$ redexes, a perfect matching in one step. Under
curried the contention is spread over the duplicator tree and the $\cstar$
channels, each one-premise, so each is a perfect matching; the enumerated
count is $n^{2}$ per erected atom. Under literal the multi-premise
constructors enumerate $n^{m}$. The enumerated-to-fired ratio \MUST{} be
reported per static channel and per scheme.
\end{hazard}

%======================================================================
\section{Addresses and freshness}\label{sec:addr}

\begin{requirement}[One allocator for both machines]\label{req:oneallocator}
As v0.2: \Comb's allocator \cite[Req.~9.5]{combspec}, factored into
\code{f1r3comb-term}.
\end{requirement}

\begin{requirement}[Deploy checks]\label{req:deploychecks}
As v0.2, citing \Comb{} v0.5: root address larger than every name of the image
(\cite[Req.~5.24]{combspec}), static channels likewise
(\cite[Req.~5.27]{combspec}), root addresses prefix-incomparable
(\cite[Req.~5.28]{combspec}), with \Comb's diagnostics. Under curried and
$\inst$ the static channels include every template's delivery channel $f$
and, under curried, every channel of the duplicator tree.
\end{requirement}

%======================================================================
\section{The operator layer}\label{sec:op}

v0.2's operator layer is unchanged in content --- species, markings,
successors with mass-action coefficients, exploration, the materialised
$T$ graded by $\Wt$, closure, $T^{\top}\chi$ --- and so is its scope: every
compiled unit with an input contains $\ev$, so triangularity and closure apply
to none of them, which \Comb{} now reports too \cite[Req.~10.2]{combspec}. One
use is added.

\begin{requirement}[Bounded multi-instance exploration]\label{req:multiexplore}
\code{explore} \MUST{} accept an instance count $k \in \{2,3\}$ and a depth
bound, deploy $k$ instances at prefix-incomparable addresses, and enumerate
every interleaving to the bound, returning the set of encoded observations at
quiescent markings. This is Requirement~\ref{req:oracle}'s oracle. The
exploration is truncated by construction and \MUST{} say so; it makes no
triangularity or closure claim.
\end{requirement}

%======================================================================
\section{Instrumentation}\label{sec:instr}

\begin{requirement}[Reports]\label{req:reports}
As v0.2, with erection attributed per scheme --- the $\cstar$ or $\inst$
firings and their $\ev$, and under curried the duplicator tree --- and with
names built per instance, sanctioned and unsanctioned mixed rows, and curried
stores detected (Hazard~\ref{haz:curriedstores}).
\end{requirement}

\begin{requirement}[The width measurement]\label{req:widthcmd}
\code{f1r3x comb-width} \MUST{} report, per program, the width distribution
under $\inst$ with erection separated from computation, and \MAY{} report the
curried scheme alongside for comparison. Under $\inst$
erection is two steps per instance; under curried it is the duplicator tree's
depth plus two. Whether either adds width without depth on real programs is
the measurement \cite[Obl.~13.7]{matnote} asks for.
\end{requirement}

\begin{requirement}[Subcommands]\label{req:fx}
As v0.2, with \code{-{}-scheme} on \code{comb-run --machine=matrix} and
\code{comb-width}, and \code{-{}-instances} on \code{comb-op}.
\end{requirement}

%======================================================================
\section{Architecture}\label{sec:arch}

The crates are v0.2's: \code{f1r3comb-term} (\Comb's, plus the rule generator
and the allocator), \code{f1r3comb-mat}, \code{f1r3comb-par},
\code{f1r3comb-op}, and the shared \code{f1r3comb-check}; none depends on
\code{f1r3comb-ir}, \code{f1r3comb-lower} or \code{f1r3comb-run}. The
test-only \code{literal-erection} feature \MUST{} live in
\code{f1r3comb-check}'s emitter and \code{f1r3comb-par}'s dev-dependencies
only.

\subsection{The device seam}\label{sec:device}

\begin{requirement}[Device kernels]\label{req:device}
As v0.2, with the erection fixed: a device backend \MUST{} target $w = 4$ and
the $\inst$ erection. Its substitution kernel is a gather over the template's
interned form with a batched insert-or-get of the results; its one erecting
join is a one-premise match at each unit's static $r$. Curried artefacts
\MUSTNOT{} reach a device --- they add a duplicator tree per unit, build more
names (E6) and cannot erect every store (E7) --- and neither may
\code{literal-erection}. A curried artefact submitted to a device backend
\MUST{} be rejected with \code{comb-mat-scheme}. It is accepted only on
identical step traces to \code{HostSequential} on the $\inst$ conformance
corpus, and no work on it \MAY{} begin before Obligation~\ref{obl:width} is
discharged. v0.2's gate on phase two's primitive is discharged by
\cite[\S6.7]{draft3}: the primitive is $\inst$.
\end{requirement}

%======================================================================
\section{Correctness obligations and testing}\label{sec:test}

\begin{obligation}[Carried from v0.2]\label{obl:carried}
Round trip; finder agreement over every generated shape, now including
$\cstar$; Python parity on E1--E4; trace certification; operator correctness;
determinism across backends, core counts and WebAssembly; the invariants of
Requirement~\ref{req:grading}; device conformance.
\end{obligation}

\begin{obligation}[The phase-two experiments, three arms]\label{obl:phase2}
\code{f1r3comb-check} \MUST{} reproduce \code{phase2-results.txt}: the E5
mixing rates for the literal scheme within a tolerance, and exactly zero for
curried and $\inst$ at every instance count under \code{MaximalProgress},
sequential choice and \code{CrossInstance}; the E6 laws exactly
(Requirement~\ref{req:dxlaw}); and E7, zero for $\inst$ and a nonzero rate
for the curried join. This is \cite[Obl.~11.3]{combspec} with the curried arm
it lacked, and the literal scheme is its negative control.
\end{obligation}

\begin{obligation}[The $W$ regression and the deciding terms]\label{obl:W}
$W$ \cite[Obl.~11.1]{combspec}, its three-instance variant, and both halves of
\cite[Req.~5.18]{combspec}, run under \code{MaximalProgress} and
\code{CrossInstance} for 100 seeds each, under each conforming scheme: the
observation at $\enc{\quo{X}}$ \MUST{} match \Bell's on every seed and no row
\MUST{} be unsanctioned-mixed. Under \code{literal-erection} the rate of the
crossed residue is recorded as a negative control.
\end{obligation}

\begin{obligation}[The marking oracle]\label{obl:marking}
Requirement~\ref{req:oracle} run on \Comb's marking corpus; any disagreement
with \Comb's analysis \MUST{} be reported to \Comb{} as a counterexample
\cite[Obl.~11.4]{combspec}.
\end{obligation}

\begin{obligation}[Addresses and differential testing]\label{obl:addr}
As v0.2, citing \cite[Obl.~11.7, 11.9, 11.12]{combspec}. \Comb{} now emits
both conforming schemes, so v0.2's test-only $\inst$ emitter is retired.
\end{obligation}

\begin{obligation}[The width of real programs]\label{obl:width}
\cite[Obl.~13.7]{matnote}, on $\inst$ artefacts, with erection separated.
Gates the device seam.
\end{obligation}

%======================================================================
\section{Errors and diagnostics}\label{sec:errors}

As v0.2, with three additions:

\begin{center}
\begin{tabularx}{\linewidth}{@{}ll>{\raggedright\arraybackslash}X@{}}
\toprule
Code & Where & Meaning\\
\midrule
\code{comb-mat-template} & import, step & A $\cstar$ template outside the
curried validity class (Requirement~\ref{req:template})\\
\code{comb-mat-curried-store} & deploy & Warning: a curried unit whose store
payload needs a join (Hazard~\ref{haz:curriedstores})\\
\code{comb-mat-scheme} & deploy & A scheme the backend does not carry: curried
or literal on a device (Requirement~\ref{req:device})\\
\bottomrule
\end{tabularx}
\end{center}

%======================================================================
\section{Performance, portability, determinism}\label{sec:perf}

As v0.2, with the targets stated for $\inst$: $10^{4}$ unfoldings of $D_{x}$
with no rise in per-unfolding step time; 64 concurrent instances of a
$10^{3}$-atom unit erected in two steps. For comparison only, the curried
scheme would take $\lceil\log_{2}(\text{atoms})\rceil + 2$.

%======================================================================
\section{What the implementation asks of the papers}\label{sec:paper}

\paragraph{Discharged since v0.2.} The multiplicity obstruction is the
paper's (\cite[Prop.~6.17, Lem.~6.19]{draft3}), and v0.2's proposal that
Remark~3.14's ``same reach'' fails under the discipline is adopted
\cite[Rem.~3.14]{draft3}; phase two has two conforming schemes
\cite[\S6.7]{draft3}; the two-stage erection is withdrawn
\cite[Rem.~5.16]{combspec}; the tag gap for constructors of constructors is
moot under either conforming scheme; the third arm \Comb{} asked for is run
(Section~\ref{sec:evidence}).

\paragraph{Open.}
\begin{enumerate}[leftmargin=1.8em]

\item \textbf{\Comb{} v0.5, Decision~5.19: make $\inst$ the default.} The
implementation path is built on $\inst$, and the machine's measurements agree
with it: equal conformance (E5), fewer steps and fewer names per unfolding
(E6), and reach over every store (E7). v0.5's three reasons for the curried
default are answered as follows. The first two --- hidden
$O(|\mathrm{unit}|)$ work and the loss of \Camp's three-line matcher --- are
properties of the content-hashed store, not of the scheme (item~4). The third
--- that $\inst$ has no home in \cite{logic}'s grammar --- is real but is a
question for the logic, not a reason to erect differently: both schemes erect
the same atoms, so the running program's occupancy is identical, and the
logic needs a production for $\inst$'s transient row only if it is to describe
erection itself (item~5). Proposed: make \code{inst-erection} the default,
keep the curried scheme behind a feature for comparison, and retire
\Mat's compatibility support once no artefact in the corpus needs it.

\item \textbf{The curried scheme's reach} (Hazard~\ref{haz:curriedstores},
E7). \cite[\S6.7]{draft3} places the address at argument positions of one
atom. A gate store whose body mentions scoped names in two atoms, or beneath a
further quotation, needs multiplicity beneath a quotation, which only
$\inst$ supplies; by \cite[Rem.~6.20]{draft3}'s own criterion the store is an
atom mentioning two marked names, so Lemma~6.19 applies. Proposed: state the
limit in \S6.7; say whether such stores are deciding, which the marking
analysis should then examine inside store payloads (\cite[Req.~5.17]{combspec});
and choose among three repairs --- widen the curried template to allow shape
formers beneath quotation, which converges on one-atom $\inst$; keep stores
static and hand the released body its parent's address so it derives the
outer names itself, which changes \cite[Req.~5.11]{combspec}; or use $\inst$
for those units.

\item \textbf{\Comb{} v0.5, Obligation~11.3.} Its $D_{x}$ figures, seven
steps and sixteen names, are $\inst$'s. Under the curried default they are ten
and twenty-two. Proposed: state the obligation per scheme.

\item \textbf{\Comb{} v0.5, Remark~5.22, first and second reasons.} The first
(``$\inst$ hides $O(|\mathrm{unit}|)$ work'') is a cost of a content-hashed
store: rebuilding a term recomputes its ancestors' hashes. With structural
interning both schemes pay in names built, and $\inst$ builds fewer (E6). The
second is a property of \code{f1r3comb-run}, not of the scheme: on this
machine both are one kernel (Decision~\ref{dec:onekernel}). Proposed: scope
both reasons to the \Camp{} machine; the third reason, the logic's grammar,
stands on its own (item~4).

\item \textbf{A place for $\inst$ in the logic} (\cite[\S14, item~4]{combspec}).
With $\inst$ on the implementation path, \cite{logic} should be read over it.
Both schemes erect the same atoms.
After erection an instance under either scheme is a multiset of base atoms
and family members, and occupancy of the running program is the same; the
schemes differ only in the transient erecting rows, whose occupancy is
uninformative in any case \cite[Req.~10.1]{combspec}. The curried template
is itself a non-name parameter, a tuple of paths and constants, so the
logic's productions would be indexed by a static parameter under either
scheme; for $\inst$ the parameter is a template, and a structural production
$\inst(\quo{\phi},\quo{\psi},\chi)$ with $\chi$ ranging over templates is the
natural addition.

\item \textbf{\Comb{} v0.5, Requirement~9.5} still describes the root address
as ``of a shape outside that program's family''; freshness is by size
(Req.~5.24, \cite[Lem.~6.14]{draft3}). Carried from v0.2.

\item \textbf{\Comb{} v0.5, Requirement~10.4} still cites \cite{logic}'s
weighted count as ``Def.~6.6''; it is Def.~4.2 and Prop.~6.17. Carried.

\item \textbf{Tags for the erecting rows and the hole.} Neither $\cstar$,
$\inst$ nor \code{HOLE} has a tag in \cite[Req.~4.11]{combspec}; the curried
template's parameter needs an encoding too. Carried and extended.

\item \textbf{\Comb{} v0.5, Requirement~9.3} (decode by projection): scope it
to builds where \code{presentation-b} is reachable. Carried.

\item \textbf{The allocator in \code{f1r3comb-term}}, and \textbf{a Merkle
form of the content hash} for $O(1)$ hashing of new addresses on
\code{f1r3comb-run}. Carried.

\item \textbf{\Comb{} v0.5, open question~5} (a bounded name universe for the
matrix encoding) was answered in v0.2: the machine multiplies the term-sized
incidence and needs no bounded universe; the closure that would is
unavailable to compiled programs for the independent reason that they
contain $\ev$.

\end{enumerate}

%======================================================================
\section{Open questions}

\begin{enumerate}[leftmargin=1.8em]
\item \textbf{Width under $\inst$} on real programs, with erection separated:
does two-step erection add width without depth, as E5 and E6 predict?
\item \textbf{Cost accounting for $\inst$}: charging by names built rather
than by step, as \cite[Rem.~5.22]{combspec} anticipates.
\item \textbf{Resolver fidelity and fairness} (\cite[Obl.~13.1,
Obl.~13.6]{matnote}).
\item \textbf{Address compaction} \cite[\S15, item~1]{combspec}.
\item \textbf{Garbage}, \textbf{the history unfolding}, \textbf{weights}: as
v0.2.
\end{enumerate}

%======================================================================
\section{Delivery}

\begin{center}
\begin{tabularx}{\linewidth}{@{}ll>{\raggedright\arraybackslash}X@{}}
\toprule
M & Deliverable & Done when\\
\midrule
X0 & \code{f1r3comb-mat} & As v0.2, with templates and \code{HOLE}\\
X1 & Finders & As v0.2, over the generated signature including the erecting
rows\\
X2 & \code{f1r3comb-par}, sequential & Resolvers; the instantiation kernel,
$\inst$ first, then the curried validity check; deploy checks; marking read
from the header; the E6 laws, $\inst$ primary\\
X3 & Phase-two evidence & The three arms of Obligation~\ref{obl:phase2}; $W$
and the deciding terms under $\inst$, and under curried for comparison;
reported to \Comb{} with the request of Section~\ref{sec:paper}, item~1\\
X4 & The marking oracle & Multi-instance bounded exploration
(Requirement~\ref{req:multiexplore}); Obligation~\ref{obl:marking} on
\Comb's corpus\\
X5 & The width gate & Under $\inst$, erection separated; needs \Comb{} M2.
\textbf{Gate} on X8\\
X6 & Threaded host & Determinism across cores and \code{wasm32}\\
X7 & \code{f1r3comb-op} & As v0.2, scoped\\
X8 & Device conformance & $\inst$ only; external backend only after
Obligation~\ref{obl:width}\\
\bottomrule
\end{tabularx}
\end{center}

%======================================================================
\section{Changes from v0.3}\label{sec:v3delta}

\begin{center}
\begin{tabularx}{\linewidth}{@{}l>{\raggedright\arraybackslash}X@{}}
\toprule
v0.3 & v0.4\\
\midrule
Curried default, following \Comb{} v0.5 & $\inst$ primary; default feature
set enables \code{inst-erection}; curried on the host for compatibility\\
Device kernels serve both conforming schemes & Device targets $\inst$ only;
curried and literal rejected with \code{comb-mat-scheme}\\
Regressions and width per scheme & $\inst$ primary, curried as comparison\\
Asks: ten & Eleven; the first asks \Comb{} to make $\inst$ its default\\
The logic's scheme left open & The logic read over $\inst$; a production for
its row proposed\\
\bottomrule
\end{tabularx}
\end{center}

\noindent
v0.3's changes from v0.2 --- the three-arm evidence, the one instantiation
kernel, the curried validity class, the marking read from the header, the
marking oracle --- carry forward unchanged.

%======================================================================
\begin{thebibliography}{99}

\bibitem{matnote} L.\,G. Meredith.
\emph{Rho combinators as sparse linear algebra.} Version~1.2,
\code{publications/rho-combinators/rhocomb-matrix}, with \code{rhocombmat.py},
\code{experiments.py} and \code{phase2.py}, September 2026.

\bibitem{matv02} F1R3FLY.io.
\emph{F1R3Comb-Mat: the rho combinators as sparse linear algebra.}
Specification v0.2, \code{publications/rho-combinators/f1r3comb-mat-spec-v02},
23 September 2026.

\bibitem{combspec} F1R3FLY.io.
\emph{F1R3Comb: a compiler from kernel rho to the name-free rho combinators,
and a machine that runs the result.} Specification v0.5,
\code{publications/rho-combinators/f1r3comb-spec-v05}, 23 September 2026.

\bibitem{draft3} L.\,G. Meredith and M. Stay.
\emph{Name-free combinators for the rho calculus: a repair, a linear encoding,
and the one operation that was missing.} Draft~3,
\code{publications/rho-combinators}, September 2026.

\bibitem{logic} F1R3FLY.io.
\emph{A spatial-behavioural logic for the rho combinators.} Draft~2,
revision~2, \code{publications/rho-combinators/rhocomb-logic}, 23 September
2026.

\bibitem{k1ndl1ng} F1R3FLY.io.
\emph{K1ndl1ng: a parser and normaliser for kernel rholang.}
\code{publications/campf1r3/k1ndl1ng-spec}, 2026.

\end{thebibliography}

\end{document}
